\noindent\unskip\SourceLeaf{19}angegeben, ein Ausdruck, der entweder mit 1, $\rho$ oder $\rho^2$ mod. $q$ congruent ist, unter $\rho$ die Einheitswurzel
$\dfrac{-1+i\sqrt{3}}{2}$ verstanden. Da nun im Gebiete der aus $\rho$ zu bildenden ganzen complexen Zahlen die Lösungen
der Congruenz $x^{q-1}\equiv 1\;(\mathrm{mod.}\ q)$ die $(q-1)$ von 0 verschiedenen rationalen ganzen Reste mod. $q$ sind,
so kommt die von \textsc{Gauss} vermöge der Congruenzen $B\equiv 0$, $B\equiv A$, $B\equiv -A\;(\mathrm{mod.}\ q)$ vollzogene Dreitheilung darauf hinaus, dass der cubische Charakter $\left[\dfrac{\pi_1}{q}\right]$ gleich 1 bez. $\rho$, $\rho^2$ ist. Im anderen Falle $q=3n+1$
ist hier der einfachste Ausdruck des Reciprocitätssatzes deshalb noch nicht erreicht, weil die Zerlegung von
$q$ in seine Primtheiler noch nicht vollzogen ist. Die Beziehung des \textsc{Gauss}'schen Theorems zum Reciprocitätsgesetze ist die folgende:

Sind die Zerlegungen der rationalen Primzahlen $p, q$ in ihre primären Primfactoren gegeben durch:
\[ p=\pi_1.\pi_2,\qquad q=\varkappa_1.\varkappa_2,\qquad \pi_k'=2\pi_k=a\pm 3b\sqrt{-3}, \]
so hat man zu setzen:
\[
\left.\begin{array}{l}
\left[\dfrac{\pi_1'}{\varkappa_1}\right]=\rho^{\varepsilon_1}\equiv {\pi_1'}^{\frac{q-1}{3}}\equiv A+B\rho-B\rho^2 \\[2mm]
\left[\dfrac{\pi_2'}{\varkappa_1}\right]=\rho^{\varepsilon_2}\equiv {\pi_2'}^{\frac{q-1}{3}}\equiv A-B\rho+B\rho^2
\end{array}\right\}(\mathrm{mod.}\ \varkappa_1).
\]

\noindent Die \textsc{Gauss}'sche Dreitheilung kommt dann mit der folgenden überein:
\[
\text{I.}\quad \varepsilon_1\equiv \varepsilon_2 (\mathrm{mod.}\ 3),\qquad \left[\dfrac{\pi_1'}{\varkappa_1}\right]=\left[\dfrac{\pi_2'}{\varkappa_1}\right],
\]
\[
\text{II.}\quad \varepsilon_1\equiv \varepsilon_2+1\;(\mathrm{mod.}\ 3),\qquad \left[\dfrac{\pi_1'}{\varkappa_1}\right]=\rho\left[\dfrac{\pi_2'}{\varkappa_1}\right],
\]
\[
\text{III.}\quad \varepsilon_1\equiv \varepsilon_2-1\;(\mathrm{mod.}\ 3),\qquad \left[\dfrac{\pi_1'}{\varkappa_1}\right]=\rho^2\left[\dfrac{\pi_2'}{\varkappa_1}\right],
\]
wobei in den drei letzten Gleichungen offenbar auch $\pi_k$ an Stelle von $\pi_k'$ treten kann. Liegt nämlich z.\,B.
der Fall II vor, so folgt durch Addition und Subtraction obiger Congruenzen in Bezug auf den Modul $\varkappa_1$:
\[
2A\equiv \rho^{\varepsilon_2}(1+\rho)\equiv -\rho^{\varepsilon_2-1},\qquad 2B\equiv -\rho^{\varepsilon_2-1}\qquad\;(\mathrm{mod.}\ \varkappa_1),
\]
so dass $(A-B)$ durch $\varkappa_1$ und als rationale Zahl demnach auch durch $q$ theilbar ist. Man hat also wirklich $B\equiv A\;(\mathrm{mod.}\ q)$. Das Reciprocitätsgesetz liefert nun:
\[
\left[\dfrac{\pi_1'}{\varkappa_1}\right]=\left[\dfrac{2}{\varkappa_1}\right].\left[\dfrac{\pi_1}{\varkappa_1}\right]=\left[\dfrac{2}{\varkappa_1}\right].\left[\dfrac{\varkappa_1}{\pi_1}\right],
\]
\[
\left[\dfrac{\pi_2'}{\varkappa_1}\right]=\left[\dfrac{2}{\varkappa_1}\right].\left[\dfrac{\pi_2}{\varkappa_1}\right]=\left[\dfrac{2}{\varkappa_1}\right].\left[\dfrac{\varkappa_1}{\pi_2}\right],
\]
so dass man für den Fall II weiter gewinnt:
\[
\left[\dfrac{\varkappa_1}{\pi_1}\right]=\rho\left[\dfrac{\varkappa_1}{\pi_2}\right]=\rho\left[\dfrac{\varkappa_2}{\pi_1}\right]^2,\qquad
\left[\dfrac{\varkappa_1}{\pi_1}\right].\left[\dfrac{\varkappa_2}{\pi_1}\right]=\left[\dfrac{q}{\pi_1}\right]=\rho,
\]
\[
\left[\dfrac{q}{\pi_1}\right]\equiv q^{\frac{p-1}{3}}\equiv -\tfrac12+\tfrac12\sqrt{-3}\equiv -\tfrac12-\tfrac12\dfrac{a}{3b}\qquad\;(\mathrm{mod.}\ \pi_1).
\]
Soll aber die rationale Zahl $q^{\frac{p-1}{3}}+\tfrac12+\tfrac12\dfrac{a}{3b}\equiv 0 \mathrm{mod.}\ \pi_1$ sein, so ist sie auch durch $\pi_2$ und also durch $p$
theilbar. Hiermit ist die unter II von \textsc{Gauss} angegebene Congruenz in Bezug auf den Modul $p$ gewonnen.
Ebenso erledigen sich die Fälle I und III.
